
Spectral properties of neural network weights are increasingly used to predict model quality without expensive evaluation. This study tested the hypothesis that low variance in randomized SVD estimates indicates well-conditioned weight matrices and better generalization. The coefficient of variation of participation ratio (CV\_PR) was computed across 100 pretrained models from the timm library using 20 random seeds per layer, then correlated with ImageNet top-1 accuracy for 94 matched models. Contrary to the hypothesis predicting r < -0.3, the observed correlation was strongly positive (Pearson r = +0.6065, p = 9.24 $\times 10^{-11})$. The 95\% confidence interval [0.506, 0.703] excludes all negative values and zero. Spearman rank correlation confirmed the result ($\rho$ = +0.6368, p = 5.26 $\times 10^{-12})$. These findings falsify the stability-as-quality assumption for this metric. Two alternative explanations are proposed: confounding by model size, and the possibility that higher spectral variance reflects richer feature representations rather than instability.
